
\section{\trKunneth{}公式}
\label{sec:kuenneth-formula}

设 $M$ 与 $N$ 是两个\trManifold，分别被赋予\trMorseFunction{} $f$ 与 $g$ 以及满足%
\trSmaleCondition{}的\trPseudogradient{}场 $X$ 与 $Y$，则 $f+g$ 是 $M\times N$ 上的
一个\trMorseFunction，而 $(X,Y)$ 是一个\trAdapted{}且满足\trSmaleCondition{}的%
\trPseudogradient{}场。$f+g$ 的诸\trCriticalPoint{}就是那些点 $(a,a')$：其中 $a$ 是 $f$
的一个\trCriticalPoint，而 $a'$ 是 $g$ 的一个\trCriticalPoint。$(a,a')$ 的\trIndex{}是
$a$ 的\trIndex{}与 $a'$ 的\trIndex{}之和。设 $(a,a')$ 是 $f+g$ 的一个\trIndex{}为 $k$ 的%
\trCriticalPoint（其中 $a$ 是 $f$ 的\trIndex{}为 $i$ 的\trCriticalPoint，而 $a'$ 是 $g$
的\trIndex{}为 $j$ 的\trCriticalPoint），并设 $(b,b')$ 是一个\trIndex{}为 $k-1$ 的%
\trCriticalPoint，使得存在 $(X,Y)$ 的把 $(a,a')$ 连接到 $(b,b')$ 的一条\trTrajectory。
$(X,Y)$ 的流为
\[
\varphi^t_{(X,Y)}(x,y)=\bigl(\varphi^t_X(x),\varphi^t_Y(y)\bigr)\text{。}
\]
特别地，
\[
\mathcal{L}_{(X,Y)}\bigl((a,a'),(b,b')\bigr)
  \cong\mathcal{L}_X(a,b)\times\mathcal{L}_Y(a',b')\text{。}
\]
若 $a\neq b$ 且 $a'\neq b'$，则要使 $\mathcal{L}_X(a,b)\times\mathcal{L}_Y(a',b')$ 非空，\footnote{译注：原书“……to be nonempty”一句中第二因子印作 $\mathcal{L}_{X'}(a',b')$；本节仅定义了场 $X$ 与 $Y$，且上行同式即作 $\mathcal{L}_Y(a',b')$，疑系排印笔误；此处径改。——译者注}
就必须有
\[
\operatorname{Ind}(a)\ge\operatorname{Ind}(b)+1
\quad\text{且}\quad
\operatorname{Ind}(a')\ge\operatorname{Ind}(b')+1\text{，}
\]
也就是说，
\[
\operatorname{Ind}(a,a')\ge\operatorname{Ind}(b,b')+2\text{。}
\]
于是，对于\trIndex{}相邻的两点 $(a,a')$ 与 $(b,b')$，我们有
\[
\mathcal{L}_{(X,Y)}\bigl((a,a'),(b,b')\bigr)=
\begin{cases}
a\times\mathcal{L}_Y(a',b') & \text{若 } a=b\\
\mathcal{L}_X(a,b)\times a' & \text{若 } a'=b'
\end{cases}
\]
从而
\[
n_{(X,Y)}\bigl((a,a'),(b,b')\bigr)=
\begin{cases}
n_Y(a',b') & \text{若 } a=b\\
n_X(a,b) & \text{若 } a'=b'\\
0 & \text{否则。}
\end{cases}
\]
映射
\[
\varPhi\colon\bigoplus_{i+j=k}C_i(f)\otimes C_j(g)\longrightarrow C_k(f+g)
\]
由
\[
\varPhi(a\otimes a')=(a,a')
\]
给出，它是一个\trAbelianGroup{}之间的\trIsomorphism。

\begin{proposition}\label{thm:kunneth-complex-isomorphism}
映射 $\varPhi$ 定义了\trComplex{}之间的一个\trIsomorphism
\[
\bigl(C_\star(f)\otimes C_\star(g),\,\partial_X\otimes 1+1\otimes\partial_Y\bigr)
  \longrightarrow\bigl(C_\star(f+g),\,\partial_{(X,Y)}\bigr)\text{。}
\]
\end{proposition}

\begin{proof}
设 $a$ 是 $f$ 的\trIndex{}为 $i$ 的一个\trCriticalPoint，并设 $a'$ 是 $g$ 的\trIndex
{}为 $j$ 的一个\trCriticalPoint。一方面，我们有
\[
\begin{aligned}
\varPhi(\partial_X\otimes 1+1\otimes\partial_Y)(a\otimes a')
  &=\varPhi\bigl(\partial_X(a)\otimes a'+a\otimes\partial_Y(a')\bigr)\\
&=\varPhi\Bigl(\sum_{b\in\operatorname{Crit}_{i-1}(f)}n_X(a,b)\,b\otimes a'\\
&\qquad+\sum_{b'\in\operatorname{Crit}_{j-1}(g)}n_Y(a',b')\,a\otimes b'\Bigr)\\
&=\sum_{b\in\operatorname{Crit}_{i-1}(f)}n_X(a,b)\,(b,a')
  +\sum_{b'\in\operatorname{Crit}_{j-1}(g)}n_Y(a',b')\,(a,b')\text{，}
\end{aligned}
\]
另一方面，我们有
\[
\begin{aligned}
\partial_{(X,Y)}\varPhi(a\otimes a')
  &=\sum_{(b,b')\in\operatorname{Crit}_{i+j-1}(f+g)}
    n_{(X,Y)}\bigl((a,a'),(b,b')\bigr)(b,b')\\
&=\sum_{b\in\operatorname{Crit}_{i-1}(f)}n_X(a,b)\,(b,a')
  +\sum_{b'\in\operatorname{Crit}_{j-1}(g)}n_Y(a',b')\,(a,b')\text{，}
\end{aligned}
\]
最后一步依据的是前面所作的注记。
\end{proof}

这个构造在 $\ZZ$ 上（考虑\trOrientation）与在 $\ZZ/2$ 上（不考虑\trOrientation）都行之
有效。\footnote{译注：原书此命题及其证明中，\trTensorProduct\trComplex{}的%
\trDifferential{}均印作 $\partial_X\otimes 1+1\otimes\partial_Y$，第二项不带任何
符号，且声称此构造在 $\ZZ$ 上（考虑\trOrientation）行之有效；对带整数系数的通常%
\trGraded{}\trTensorProduct\trComplex，第二项疑应带符号 $(-1)^i$（即 Koszul 符号
约定），否则 $d^2=0$ 一般不成立；此处照原书排印，存疑备案。——译者注}
不过，等到取%
\trHomology{}时，在 $\ZZ/2$ 上工作更为便利（所需的\trHomology{}代数可在\cref{subsec:appB-kunneth-z2}中
找到）。
由此我们导出下述结果。

\begin{corollary}\label{thm:mod2-tensor-homology}
积 $M\times N$ 上的\trComplex{} $(C_\star(f+g),\partial_{(X,Y)})$ 的模 2 \trHomology{}%
\trIsomorphic{}于 $(C_\star(f),\partial_X)$ 与 $(C_\star(g),\partial_Y)$ 的模 2 %
\trHomology{}的\textbf{\trTensorProduct}（tensor product）。
\end{corollary}

\begin{corollary}[\trKunneth{}公式]\label{thm:kunneth-formula}\BookIndexDefinition{trKunnethFormula}
若 $M$ 与 $N$ 是两个\trCompact{}\trManifold，则存在一个\trIsomorphism
\[
HM_k(M\times N;\ZZ/2)\longrightarrow
\bigoplus_{i+j=k}HM_i(M;\ZZ/2)\otimes HM_j(N;\ZZ/2)\text{。}
\]
\end{corollary}

证明是\cref{thm:kunneth-complex-isomorphism}与\cref{thm:appB-kunneth-tensor-product}的直接推论。
\noqed
